\emph{Premise.} Hypothesis (P) of the statement is Premise \claimref{synthesis.premise-positivity-of-values}; every statement below about the values of states is under it, relative to the fixed split of the generators.

\emph{Preliminaries.} By Theorem \claimref{synthesis.signed-directions-and-the-integer-clifford-algebra} the algebra over $\Z$ is free on the blades $e_{S}$, the images of the group $G$ of signed directions, with $e_{S}e_{T} = \pm e_{S \triangle T}$; by Theorem \claimref{synthesis.division-gives-the-rationals} the same holds with coefficients in $\Q$, and by base change in any field $k$ of characteristic $0$. Inversion is an anti-automorphism of $G$ fixing the sign bit, so its linear extension $\dagger$ is an anti-automorphism with $x^{\dagger\dagger} = x$. The product $e_{S}^{\dagger}e_{T} = \pm e_{S \triangle T}$ is a scalar only for $S = T$, where it is $1$; hence for $x = \sum x_{S}e_{S}$, $y = \sum y_{S}e_{S}$,
\[
\langle x^{\dagger}y\rangle_{0} = \sum_{S}x_{S}y_{S}, \qquad \langle x^{\dagger}x\rangle_{0} = \sum_{S}x_{S}^{2} \ge 0, \qquad \langle xy\rangle_{0} = \langle yx\rangle_{0},
\]
the last because $e_{S}e_{T}$ and $e_{T}e_{S}$ are scalars only for $S = T$, where they agree. In an ordered field a sum of squares is non-negative and vanishes only if every term does. For a state $\omega$ and a self-adjoint $h$, $0 \le \omega((h - t)(h - t)^{\dagger}) = \omega(h^{2}) - 2t\omega(h) + t^{2}$ for every scalar $t$; at $t = \omega(h)$ this gives $\omega(h)^{2} \le \omega(h^{2})$. For any $u, v$ and scalar $t$, $0 \le \omega((u + tv)^{\dagger}(u + tv))$ is a quadratic polynomial in $t$, so
\[
\bigl(\omega(u^{\dagger}v) + \omega(v^{\dagger}u)\bigr)^{2} \le 4\,\omega(u^{\dagger}u)\,\omega(v^{\dagger}v). \tag{CS}
\]
(Here and below $\omega(x^{\dagger}x) \ge 0$ is used in place of $\omega(xx^{\dagger}) \ge 0$; the two are the same condition, applied to $x^{\dagger}$.) For self-adjoint $\rho$, $\langle\rho x^{\dagger}\rangle_{0} = \langle(x\rho)^{\dagger}\rangle_{0} = \langle x\rho\rangle_{0} = \langle\rho x\rangle_{0}$, since $\langle y^{\dagger}\rangle_{0} = \langle y\rangle_{0}$; so every functional $x \mapsto \langle\rho x\rangle_{0}$ with $\rho^{\dagger} = \rho$ has the reality property of (P). The frame remark of the statement: in signature $(p,0)$ the adjoint is the reverse and in $(0,q)$ the Clifford conjugate, both defined without reference to a frame; in indefinite signature clause (vi) gives two orthonormal frames with different adjoints.

\emph{(i).} $(A_{i}B_{j})^{\dagger} = B_{j}A_{i} = A_{i}B_{j}$, so $C^{\dagger} = C$. Put $S = B_{1} + B_{2}$, $D = B_{1} - B_{2}$; they commute with $A_{1}, A_{2}$, and $S^{2} + D^{2} = 4$, $SD = -Y$, $DS = Y$. So $C^{2} = (A_{1}S + A_{2}D)^{2} = S^{2} + D^{2} + A_{1}A_{2}SD + A_{2}A_{1}DS = 4 - (A_{1}A_{2} - A_{2}A_{1})Y = 4 - XY$.

\emph{(ii).} $X^{\dagger} = A_{2}A_{1} - A_{1}A_{2} = -X$ and likewise $Y^{\dagger} = -Y$, while $S_{A}^{\dagger} = S_{A}$, $S_{B}^{\dagger} = S_{B}$. Since $A_{1}^{2} = A_{2}^{2} = 1$, $X^{2} = A_{1}A_{2}A_{1}A_{2} + A_{2}A_{1}A_{2}A_{1} - 2$ and $S_{A}^{2} = A_{1}A_{2}A_{1}A_{2} + A_{2}A_{1}A_{2}A_{1} + 2$, so
\[
X^{\dagger}X = -X^{2} = 4 - S_{A}^{2}, \qquad Y^{\dagger}Y = 4 - S_{B}^{2}.
\]
$X$ and $Y$ commute, so $Z^{\dagger}Z = (Y - X)(X - Y) = -X^{2} - Y^{2} + 2XY = 8 - S_{A}^{2} - S_{B}^{2} + 2XY$. Therefore $\tfrac12(Z^{\dagger}Z + S_{A}^{2} + S_{B}^{2}) = 4 + XY = 8 - C^{2}$.

\emph{(iii).} Positivity of $\langle x^{\dagger}x\rangle_{0}$ is in the preliminaries. For $\rho = \sum_{r}y_{r}y_{r}^{\dagger}$, $\langle\rho\rangle_{0} = \sum_{r}\langle y_{r}^{\dagger}y_{r}\rangle_{0} > 0$ unless every $y_{r} = 0$, and $\langle\rho\,xx^{\dagger}\rangle_{0} = \sum_{r}\langle y_{r}^{\dagger}xx^{\dagger}y_{r}\rangle_{0} = \sum_{r}\langle(x^{\dagger}y_{r})^{\dagger}(x^{\dagger}y_{r})\rangle_{0} \ge 0$ by cyclicity; so $\omega_{\rho}$ is non-negative on every $xx^{\dagger}$; $\rho$ is self-adjoint, so $\omega_{\rho}$ is real, and it is a state. For any state, $\omega(C)^{2} \le \omega(C^{2})$ since $C$ is self-adjoint, and applying $\omega$ to the certificate of (ii) gives the displayed equality; the three terms are values of $Z^{\dagger}Z$, $S_{A}S_{A}^{\dagger}$, $S_{B}S_{B}^{\dagger}$, hence non-negative, which gives $\omega(C^{2}) \le 8$. If $\omega(C)^{2} = 8$ then $\omega(C^{2}) = 8$ and the three non-negative terms sum to zero, so each is zero; and $2\,(\omega(C)/2)^{2} = 4$ shows that $2$ is a square in $k$, which excludes $k = \Q$. For the second inequality take $u = X$, $v = Y$ in (CS): $u^{\dagger}v = -XY = v^{\dagger}u$, so $4\,\omega(XY)^{2} \le 4\,\omega(X^{\dagger}X)\,\omega(Y^{\dagger}Y) = 4\,\omega(4 - S_{A}^{2})\,\omega(4 - S_{B}^{2})$, and $\omega(C^{2}) - 4 = -\omega(XY)$. If $A_{1}$ and $A_{2}$ commute then $X = 0$ and $C^{2} = 4$; likewise for $B$.

\emph{(iv).} A blade is invertible with inverse $\pm$ itself, so a self-adjoint blade has square $1$, and two blades commute or anticommute. $X = 2A_{1}A_{2}$ and $Y = 2B_{1}B_{2}$, so by (i) $C^{2} = 4 - 4A_{1}A_{2}B_{1}B_{2} = 4 + 4K$. $K$ is a blade up to sign; $K^{\dagger} = -B_{2}B_{1}A_{2}A_{1} = -A_{1}A_{2}B_{1}B_{2} = K$ (two sign changes), so $K^{2} = KK^{\dagger} = 1$. $A_{1}$ anticommutes with $A_{1}A_{2}$ and commutes with $B_{1}B_{2}$, and likewise for the other three, so each $A_{i}B_{j}$ commutes with $K$ and $KC = CK$. $K$ is not a scalar: otherwise $A_{1}A_{2} = \pm B_{1}B_{2}$ would commute with $A_{1}$. Since the $A$'s commute with the $B$'s, $KA_{1}B_{1} = -(A_{1}A_{2}A_{1})(B_{1}B_{2}B_{1}) = -(-A_{2})(-B_{2}) = -A_{2}B_{2}$, and in the same way $KA_{2}B_{2} = -A_{1}B_{1}$, $KA_{1}B_{2} = A_{2}B_{1}$, $KA_{2}B_{1} = A_{1}B_{2}$; so $KC = C$ and $C^{3} = 4C + 4KC = 8C$. Each $A_{i}B_{j}$ is a blade that is not a scalar ($A_{i} = \pm B_{j}$ is excluded because $B_{j}$ commutes with both $A$'s), so $\langle C\rangle_{0} = \langle KC\rangle_{0} = 0$, and $\langle K\rangle_{0} = 0$. \emph{The functionals.} $\rho_{r} = (1 + K)(1 + rC)$ is self-adjoint, so $\omega_{r}$ is real. $\omega_{r}(1) = \langle 1 + K + rC + rKC\rangle_{0} = 1$; $\omega_{r}(C) = \langle(1 + K)C\rangle_{0} + r\langle(1 + K)(4 + 4K)\rangle_{0} = 8r$; $\omega_{r}(C^{2}) = \langle(1 + K)(4 + 4K)\rangle_{0} + r\langle(1 + K)C^{3}\rangle_{0} = 8 + 8r\langle(1 + K)C\rangle_{0} = 8$. Positivity: put $\Pi = \tfrac12(1 + K)$, a self-adjoint idempotent commuting with $C$, and $y = \Pi x$. By cyclicity, $\tfrac12\omega_{r}(xx^{\dagger}) = \langle x^{\dagger}\Pi(1 + rC)\Pi x\rangle_{0}$, which is $\langle y^{\dagger}y\rangle_{0} + r\langle y^{\dagger}Cy\rangle_{0}$. By the Cauchy-Schwarz inequality for the positive symmetric form $\langle a^{\dagger}b\rangle_{0}$ (an inequality between squares, valid in any ordered field), $\langle y^{\dagger}Cy\rangle_{0}^{2} \le \langle y^{\dagger}y\rangle_{0}\langle y^{\dagger}C^{2}y\rangle_{0} = 8\langle y^{\dagger}y\rangle_{0}^{2}$, because $C^{2}\Pi = 8\Pi$. So $(r\langle y^{\dagger}Cy\rangle_{0})^{2} \le 8r^{2}\langle y^{\dagger}y\rangle_{0}^{2} \le \langle y^{\dagger}y\rangle_{0}^{2}$ when $8r^{2} \le 1$, and $\omega_{r}(xx^{\dagger}) \ge 0$. Conversely if $\omega_{r}$ satisfies (P) then by (iii) $64r^{2} = \omega_{r}(C)^{2} \le \omega_{r}(C^{2}) = 8$. For rational $v$ with $v^{2} < 8$, $r = v/8$ has $8r^{2} < 1$. The rational $r$ with $8r^{2} \le 1$ have supremum $\sqrt{2}/4$, which is not rational, and by (iii) no state has $\omega(C)^{2} = 8$ over $\Q$; a field containing $\sqrt{2}$ contains $r = \sqrt{2}/4$.

\emph{(v).} $1 + C$ is self-adjoint, so $\lambda$ is real; $\lambda(1) = 1 + \langle C\rangle_{0} = 1$ and $\lambda(C) = \langle C\rangle_{0} + \langle C^{2}\rangle_{0} = \langle 4 + 4K\rangle_{0} = 4$ by (iv). $\lambda((C - t)^{2}) = \langle(1 + C)(C^{2} - 2tC + t^{2})\rangle_{0} = \langle C^{2}\rangle_{0} + t^{2} + \langle C^{3}\rangle_{0} - 2t\langle C^{2}\rangle_{0} = 4 + t^{2} - 8t$, using $\langle C\rangle_{0} = 0$ and $\langle C^{3}\rangle_{0} = 8\langle C\rangle_{0} = 0$; at $t = 1$ it is $-3$. $C - t$ is self-adjoint, so this is a value on an element $xx^{\dagger}$.

\emph{(vi).} $e_{1}$ and $f_{1}$ anticommute, so $A^{2} = (\tfrac{17}{8})^{2} - (\tfrac{15}{8})^{2} = \tfrac{289 - 225}{64} = 1$, and $A^{\dagger} = \tfrac{17}{8}e_{1} - \tfrac{15}{8}f_{1} \ne A$. $p_{2} = f_{2}e_{2}$ is an even element of the other mode, so it commutes with $e_{1}$, $f_{1}$; it is self-adjoint with square $1$. $e_{1}$ and $p_{2}$ are commuting observables, so $\rho = \tfrac12(1 + e_{1})\,\tfrac12(1 + p_{2})$ is a self-adjoint idempotent, $\rho = \rho\rho^{\dagger}$, and $\omega_{\rho}$ is a state with $\langle\rho\rangle_{0} = \tfrac14$. $C = 2Ap_{2}$, and $\langle\rho\,e_{1}p_{2}\rangle_{0} = \tfrac14$ (the term $e_{1}p_{2}$ of $\rho$), $\langle\rho\,f_{1}p_{2}\rangle_{0} = 0$; so $\omega_{\rho}(C) = 2 \cdot \tfrac{17}{8} = \tfrac{17}{4}$. $C^{2} = 4A^{2}p_{2}^{2} = 4$. The step $\omega(C)^{2} \le \omega(C^{2})$ fails because $C$ is not self-adjoint. \emph{The boosted frame.} $F = \tfrac{15}{8}e_{1} + \tfrac{17}{8}f_{1}$ has $F^{2} = \tfrac{225 - 289}{64} = -1$ and $AF + FA = 2 \cdot \tfrac{17 \cdot 15}{64}(e_{1}^{2} + f_{1}^{2}) = 0$, so $A, F, e_{2}, f_{2}$ are orthonormal generators of the same algebra with the same signs, and the adjoint $\dagger'$ defined from them fixes $A$. Then $(1 - A)(1 - A)^{\dagger'} = (1 - A)^{2} = 2 - 2A$, and $\omega_{\rho}(2 - 2A) = 2 - 2 \cdot \tfrac{17}{8} = -\tfrac94$.

\emph{(vii).} The adjoint fixes $e_{j}$ (square $1$) and negates $f_{j}$ (square $-1$), and $p_{j}^{\dagger} = e_{j}^{\dagger}f_{j}^{\dagger} = -e_{j}f_{j} = f_{j}e_{j} = p_{j}$, with $p_{j}^{2} = -f_{j}^{2}e_{j}^{2} = 1$ and $e_{j}p_{j} = -p_{j}e_{j}$. The sub-Clifford algebra of mode $j$ has basis $1, e_{j}, f_{j}, p_{j}$, so its self-adjoint part is spanned by $1, e_{j}, p_{j}$, and $(s + xe_{j} + zp_{j})^{2} = s^{2} + x^{2} + z^{2} + 2s(xe_{j} + zp_{j})$ equals $1$ exactly when $s = 0$, $x^{2} + z^{2} = 1$ or $x = z = 0$, $s = \pm 1$. $e_{1}e_{2} = -e_{2}e_{1}$ has square $-1$ and adjoint $e_{2}e_{1} = -e_{1}e_{2}$. The twist is that of Theorem \claimref{tower.commuting-copies} with one plane (mode $1$) and ground algebra mode $2$; the lines needed are proved here. $p_{1}$ is an even element of mode $1$, so it commutes with $e_{2}$, $f_{2}$, and it anticommutes with $e_{1}$, $f_{1}$, with $p_{1}^{2} = 1$. Hence for $y, y' \in \{e_{2}, f_{2}\}$: $(yp_{1})^{2} = y^{2}p_{1}^{2} = y^{2}$; $(yp_{1})(y'p_{1}) = yy'$, so $e_{2}p_{1}$ and $f_{2}p_{1}$ anticommute; and for $g \in \{e_{1}, f_{1}\}$, $g\,yp_{1} = -y\,gp_{1} = y\,p_{1}g$, so $yp_{1}$ commutes with mode $1$. The twisted generators generate a copy $B'$ of $\Cl(1,1)$; its parity is $(f_{2}p_{1})(e_{2}p_{1}) = p_{2}$ and $(e_{2}p_{1})^{\dagger} = p_{1}e_{2} = e_{2}p_{1}$, so the self-adjoint part of $B'$ is spanned by $1$, $e_{2}p_{1}$, $p_{2}$ and its directions are $x\,e_{2}p_{1} + z\,p_{2}$. For directions $a_{1}, a_{2}$ of $A$: $a_{1}a_{2} = a_{1} \cdot a_{2} + (a_{1} \wedge a_{2})\,e_{1}p_{1}$, so $S_{A} = 2a_{1} \cdot a_{2}$ and $X = 2(a_{1} \wedge a_{2})e_{1}p_{1} = -2(a_{1} \wedge a_{2})f_{1}$; likewise $Y = 2(b_{1} \wedge b_{2})\,e_{2}p_{1}p_{2} = -2(b_{1} \wedge b_{2})f_{2}p_{1}$. Then $f_{1}f_{2}p_{1} = f_{1}f_{2}f_{1}e_{1} = f_{2}e_{1} = -e_{1}f_{2}$, so $XY = -4(a_{1} \wedge a_{2})(b_{1} \wedge b_{2})K$ and $C^{2} = 4 - XY$ is as stated. $K^{\dagger} = f_{2}^{\dagger}e_{1}^{\dagger} = -f_{2}e_{1} = K$ and $K^{2} = -e_{1}^{2}f_{2}^{2} = 1$. $K = e_{1}f_{2}$ anticommutes with $e_{1}$, $p_{1}$, $e_{2}p_{1}$ and $p_{2}$, so it commutes with every product of a direction of $A$ and a direction of $B$, whatever the wedges. Since $K$ is self-adjoint with square $1$, $\omega(K)^{2} \le 1$, so $\omega(C)^{2} \le \omega(C^{2}) = 4 + 4(a_{1} \wedge a_{2})(b_{1} \wedge b_{2})\omega(K) \le 4 + 4|a_{1} \wedge a_{2}||b_{1} \wedge b_{2}|$. With $(a_{1} \wedge a_{2})^{2} = 1 - (a_{1} \cdot a_{2})^{2} \le 1$, the value $8$ requires both wedges to have square $1$, that is $a_{1} \cdot a_{2} = b_{1} \cdot b_{2} = 0$, which is $S_{A} = S_{B} = 0$. \emph{Parity.} $\Gamma = p_{1}p_{2}$ anticommutes with $e_{1}$ and with $e_{2}p_{1}$ and commutes with $p_{1}$, $p_{2}$; so $\tfrac12(x + \Gamma x\Gamma)$ removes the $e_{1}$ and $e_{2}p_{1}$ components of the directions, the averaged directions are $z_{i}p_{1}$ and $z'_{j}p_{2}$, and their CHSH combination is the stated multiple of $p_{1}p_{2}$, with $|z_{1}(z'_{1} + z'_{2}) + z_{2}(z'_{1} - z'_{2})| \le |z'_{1} + z'_{2}| + |z'_{1} - z'_{2}| \le 2$ since all $|z| \le 1$. Theorem \claimref{singularity.descent-parity-average} is stated for one mode; its total-parity form, used here, is a definition.

\emph{(viii).} Expanding with $p_{1}e_{2}p_{1} = e_{2}$ and $e_{1}e_{2}p_{1} = e_{1}e_{2}f_{1}e_{1} = e_{2}f_{1}$ gives $C$; the wedges are $1$, so $C^{2} = 4(1 + K)$. $K$ anticommutes with each of the four directions, so $KA_{1}B_{1} = -(A_{1}A_{2}A_{1})(B_{1}B_{2}B_{1})$ with $K = -A_{1}A_{2}B_{1}B_{2}$ (from $XY = 4A_{1}A_{2}B_{1}B_{2}$), which is $-A_{2}B_{2}$; in the same way $KA_{2}B_{2} = -A_{1}B_{1}$, $KA_{1}B_{2} = A_{2}B_{1}$, $KA_{2}B_{1} = A_{1}B_{2}$. So $KC = C$ and $C^{3} = 4C + 4KC = 8C$. Then $P^{2} = \tfrac{1}{16}(1 + K)^{2} + \tfrac{\sqrt{2}}{16}(1 + K)C + \tfrac{2}{64}C^{2} = \tfrac18(1 + K) + \tfrac{\sqrt{2}}{8}C + \tfrac18(1 + K) = P$, $P^{\dagger} = P$, and $CP = \tfrac14(C + KC) + \tfrac{\sqrt{2}}{8}C^{2} = \tfrac12C + \tfrac{\sqrt{2}}{2}(1 + K) = 2\sqrt{2}\,P$. The four terms of $C$ and $K$ are blades other than $1$, so $\langle P\rangle_{0} = \tfrac14$. The identity $PxP = 4\langle Px\rangle_{0}P$ is linear in $x$ and is verified on the sixteen blades by the check (with coefficients $a + b\sqrt{2}$, $a, b \in \Q$, as exact pairs). $\omega_{P}$ is a state because $P = PP^{\dagger}$ is self-adjoint; $4P = (1 + K)(1 + \tfrac{\sqrt{2}}{4}C)$ because $KC = C$, so $\omega_{P}$ is the $\omega_{r}$ of (iv) at $r = \sqrt{2}/4$; and $\omega_{P}(C) = 4\langle PC\rangle_{0} = 4 \cdot 2\sqrt{2}\langle P\rangle_{0} = 2\sqrt{2}$. The blades of $P$ in mode $1$ alone are $\tfrac14$; those in mode $2$ alone are $\tfrac14$ and the term $e_{2}$ of $C$. $p_{1}p_{2}$ commutes with $K$, $e_{2}f_{1}$ and $p_{1}p_{2}$ and anticommutes with the odd terms $e_{1}p_{2}$ and $e_{2}$ of $C$.

\emph{Uniqueness on self-adjoint elements.} Let $\omega$ be linear with $\omega(1) = 1$, $\omega(xx^{\dagger}) \ge 0$ and $\omega(C) = 2\sqrt{2}$ (reality is not used until the last line) and $D = C - 2\sqrt{2}$. By (iii), $\omega(D^{2}) = \omega(C^{2}) - 8 \le 0$, so $\omega(D^{2}) = 0$. \emph{Bounded elements.} For $y = \sum_{S}y_{S}e_{S}$ and $c = (\sum_{S}|y_{S}|)^{2}$,
\[
c - y^{\dagger}y = \tfrac12\sum_{S,T}|y_{S}y_{T}|\,(e_{S} - \epsilon_{ST}e_{T})^{\dagger}(e_{S} - \epsilon_{ST}e_{T}), \qquad \epsilon_{ST} = \sgn(y_{S}y_{T}),
\]
since $(e_{S} - \epsilon e_{T})^{\dagger}(e_{S} - \epsilon e_{T}) = 2 - \epsilon(e_{S}^{\dagger}e_{T} + e_{T}^{\dagger}e_{S})$; so $D(c - y^{\dagger}y)D$ is a sum of terms $t\,w^{\dagger}w$ with $t \ge 0$, and $\omega((yD)^{\dagger}(yD)) \le c\,\omega(D^{2}) = 0$. By (CS) with $u = x$, $v = yD$: $\omega(x^{\dagger}yD) + \omega(Dy^{\dagger}x) = 0$ for all $x, y$. Now $1 - P = -\tfrac{1}{16}(C - 2\sqrt{2})(C + 4\sqrt{2}) = gD = Dg$ with $g = -\tfrac{1}{16}(C + 4\sqrt{2})$ self-adjoint (the product is $-\tfrac{1}{16}(C^{2} + 2\sqrt{2}C - 16)$, and $C^{2} = 4(1 + K)$). For self-adjoint $h$, $h - PhP = (1 - P)h + Ph(1 - P) = Dgh + PhgD$. With $x = 1$, $y = hg$: $\omega(hgD) + \omega(Dgh) = 0$. So $\omega(h - PhP) = -\omega(hgD) + \omega(PhgD) = -\omega((1 - P)hgD) = -\omega(DmD)$ with $m = ghg$ self-adjoint. Finally $DmD = \tfrac14\bigl(((1 + m)D)^{\dagger}(1 + m)D - ((1 - m)D)^{\dagger}(1 - m)D\bigr)$, and both terms have value $0$ by the bound just proved. Hence $\omega(h) = \omega(PhP) = 4\langle Ph\rangle_{0}\,\omega(P)$, and $h = 1$ gives $\omega(P) = 1$: $\omega(h) = \omega_{P}(h)$. A state with $\omega(x^{\dagger}) = \omega(x)$ is determined by its values on self-adjoint elements. So under (P) it equals $\omega_{P}$.

\emph{(ix).} $p_{3}$ is self-adjoint of square $1$ and commutes with the generators of modes $1$ and $2$, so $P_{\pm}$ are self-adjoint idempotents with $P_{+}P_{-} = 0$, commuting with $A_{1}$, $e_{1}$, $B_{1}$, $B_{2}$. Then $A_{2}$ is self-adjoint, $A_{2}^{2} = e_{1}^{2}P_{+} + p_{1}^{2}P_{-} = 1$, it commutes with $B_{1}, B_{2}$ (as $e_{1}$ and $p_{1}$ do, by (vii)), and $S_{A} = (p_{1}e_{1} + e_{1}p_{1})P_{+} + 2P_{-} = 1 - p_{3}$, while $X = [A_{1},A_{2}] = 2p_{1}e_{1}P_{+}$. By (i), $C^{2} = 4 - 2p_{1}e_{1}P_{+} \cdot 2B_{1}B_{2} = 4 + 4K'P_{+}$. Let $C_{0}$ be the CHSH element of the blade configuration $(p_{1}, e_{1}; p_{2}, e_{2}p_{1})$; $K'$ is its element $K$ of (iv), so $K'^{\dagger} = K'$, $K'^{2} = 1$, $K'C_{0} = C_{0}K' = C_{0}$ and $C_{0}^{2} = 4 + 4K'$. Since $A_{2}P_{+} = e_{1}P_{+}$, $CP_{+} = C_{0}P_{+}$. Put $\Pi = \tfrac14(1 + K')(1 + p_{3})$, a self-adjoint idempotent (a product of commuting ones) with $\Pi = \Pi P_{+}$ and $\Pi K' = \Pi$. Then $C\Pi = C_{0}\Pi = \Pi C_{0} = \Pi C$, $C^{2}\Pi = (4 + 4K')\Pi = 8\Pi$ and $S_{A}\Pi = (1 - p_{3})\Pi = 0$. The functional is $\omega'_{r}(x) = 4\langle\Pi(1 + rC)x\rangle_{0}$, real because $\Pi(1 + rC)$ is self-adjoint. The blades of $C_{0}$, $K'$ and $K'C_{0} = C_{0}$ are not scalars and do not contain the generators of mode $3$, so they stay non-scalar after multiplication by $p_{3}$: $4\langle\Pi\rangle_{0} = 1$ and $4\langle\Pi C\rangle_{0} = 4\langle\Pi C_{0}\rangle_{0} = 0$. Hence $\omega'_{r}(1) = 1$, $\omega'_{r}(C) = 4r\langle\Pi C^{2}\rangle_{0} = 8r$, $\omega'_{r}(C^{2}) = 8 + 8r \cdot 4\langle\Pi C\rangle_{0} = 8$, and $\omega'_{r}(S_{A}^{2}) = 4\langle(1 + rC)S_{A}^{2}\Pi\rangle_{0} = 0$. Positivity is as in (iv): with $y = \Pi x$, $\tfrac14\omega'_{r}(xx^{\dagger}) = \langle y^{\dagger}y\rangle_{0} + r\langle y^{\dagger}Cy\rangle_{0}$ and $\langle y^{\dagger}Cy\rangle_{0}^{2} \le \langle y^{\dagger}y\rangle_{0}\langle y^{\dagger}C^{2}y\rangle_{0} = 8\langle y^{\dagger}y\rangle_{0}^{2}$, so $8r^{2} \le 1$ gives $\omega'_{r}(xx^{\dagger}) \ge 0$.
